\documentclass[11pt,a4paper]{article}
\usepackage[T1]{fontenc}
\usepackage[utf8]{inputenc}
\usepackage{amsmath,amssymb,xcolor}
\usepackage[margin=25mm]{geometry}
\usepackage[colorlinks=true,urlcolor=blue]{hyperref}
\definecolor{nexorablue}{HTML}{2459B5}
\setlength{\parindent}{0pt}
\setlength{\parskip}{7pt}
\setlength{\emergencystretch}{2em}
\renewcommand{\familydefault}{\sfdefault}
\begin{document}

{\small\bfseries\color{nexorablue} NEXORA CPGE}\par
{\small MP / NEX-RED-03 / Corrigé}\par
{\LARGE\bfseries Une matrice triangulaire sans base propre}\par
\vspace{0.25em}\hrule\vspace{0.6em}

\section*{Dimensions}
On a $\chi_A=(X-4)^2(X+1)$. Dans $\ker(A-4I)$, la seconde et la troisième coordonnées sont nulles. Dans $\ker(A+I)$, seules les troisièmes coordonnées sont libres. On obtient deux droites propres, insuffisantes en dimension trois.

\section*{Polynôme minimal}
Le premier bloc exige le facteur $(X-4)^2$ et le dernier bloc exige $X+1$. Ces facteurs sont premiers entre eux. Leur produit, qui annule les deux blocs, est le polynôme minimal.


\end{document}
